% !PRINTED 200
\weberpage{200}
%<<CONT
\[
\begin{aligned}
& A = \tfrac{1}{2} a_0^2 (U^2 + V^2 + W^2) \\
(7)\quad & B = a_0^3 \bigl[U^2(V - W) + V^2(W - U) + W^2(U - V)\bigr] \\
& \phantom{B} = a_0^3 (W - V)(U - W)(V - U).
\end{aligned}
\]
Zur Berechnung dieser Grössen sind die nöthigen Formeln
schon in \S.\,47 entwickelt. Nach den dortigen Formeln (10) ergiebt sich
\[
U = v^2 - w^2, \quad V = w^2 - u^2, \quad W = u^2 - v^2,
\]
wenn $u^2$, $v^2$, $w^2$ die Wurzeln der cubischen Resolvente
\[
z^3 + 2 a z^2 + (a^2 - 4 c) z - b^2 = 0
\]
sind, und es folgt daraus
\[
\begin{aligned}
A &= a_0^2 \bigl[(u^2 + v^2 + w^2)^2 - 3 (v^2 w^2 + w^2 u^2 + u^2 v^2)\bigr] \\
&= a_0^2 (a^2 + 12 c) \\
B &= a_0^3 (2 u^2 - v^2 - w^2)(2 v^2 - w^2 - u^2)(2 w^2 - u^2 - v^2) \\
&= a_0^3 (3 u^2 + 2 a)(3 v^2 + 2 a)(3 w^2 + 2 a) \\
&= a_0^3 (2 a^3 - 72 a c + 27 b^2),
\end{aligned}
\]
so dass $A$, $B$ dieselbe Bedeutung haben, wie in \S.\,47 und durch
die Coëfficienten $a_0, a_1, a_2, a_3, a_4$ ausgedrückt, so dargestellt sind:
\weq{8}{A = a_2^2 - 3 a_1 a_3 + 12 a_0 a_4}
\weq{9}{B = 27 a_1^2 a_4 + 27 a_0 a_3^2 + 2 a_2^3 - 72 a_0 a_2 a_4 - 9 a_1 a_2 a_3.}
Auch die Discriminante der biquadratischen Gleichung als
eine dritte, aber von $A$ und $B$ abhängige Invariante
\weq{10}{D = a_0^6 \, U^2 V^2 W^2}
haben wir an der erwähnten Stelle schon gebildet und gefunden
\weq{11}{27 D = 4 A^3 - B^2.}
% [sic] Der Druck verweist auf Paragraph 67; die hier wörtlich angeführte Formel ist die binäre Spezialisierung von Paragraph 60, (3), während Paragraph 67 erst später beginnt.
Nach der Formel des \S.\,67, $2\lambda = n\mu - \nu$, erhalten wir für die
bis jetzt gefundenen invarianten Bildungen, die wir für eine
transformirte Form mit Accenten bezeichnen, folgende Relationen
\weq{12}{
\begin{aligned}
& H' = r^2 H, \quad T' = r^3 T, \\
& A' = r^4 A, \quad B' = r^6 B, \quad D' = r^{12} D.
\end{aligned}
}
